\documentclass[11pt,fontset=fandol]{ctexart}
\usepackage[a4paper,margin=23mm]{geometry}
\usepackage{amsmath,amssymb,mathtools,bm,mathrsfs,esint}
\usepackage{graphicx,longtable,booktabs,array,calc}
\usepackage[normalem]{ulem}
\usepackage{xcolor,hyperref}
\hypersetup{hidelinks}
\setcounter{secnumdepth}{-1}
\newcounter{none}
\providecommand{\tightlist}{\setlength{\itemsep}{0pt}\setlength{\parskip}{0pt}}
\providecommand{\passthrough}[1]{#1}
\usepackage{amssymb,mathrsfs,esint,bm,newunicodechar}
\allowdisplaybreaks
\setlength{\emergencystretch}{3em}
\providecommand{\llangle}{\langle\!\langle}
\providecommand{\rrangle}{\rangle\!\rangle}
\newunicodechar{∞}{\ensuremath{\infty}}
\newunicodechar{∈}{\ensuremath{\in}}
\newunicodechar{∑}{\ensuremath{\sum}}
\newunicodechar{∫}{\ensuremath{\int}}
\newunicodechar{≤}{\ensuremath{\leq}}
\newunicodechar{≥}{\ensuremath{\geq}}
\newunicodechar{≠}{\ensuremath{\neq}}
\newunicodechar{→}{\ensuremath{\to}}
\newunicodechar{⇒}{\ensuremath{\Rightarrow}}
\newunicodechar{α}{\ensuremath{\alpha}}
\newunicodechar{β}{\ensuremath{\beta}}
\newunicodechar{γ}{\ensuremath{\gamma}}
\newunicodechar{λ}{\ensuremath{\lambda}}
\newunicodechar{φ}{\ensuremath{\varphi}}
\newunicodechar{ε}{\ensuremath{\varepsilon}}
\newunicodechar{σ}{\ensuremath{\sigma}}
\newunicodechar{θ}{\ensuremath{\theta}}
\newunicodechar{δ}{\ensuremath{\delta}}
\newunicodechar{η}{\ensuremath{\eta}}
\newunicodechar{π}{\ensuremath{\pi}}
\newunicodechar{ξ}{\ensuremath{\xi}}
\newunicodechar{Σ}{\ensuremath{\Sigma}}
\newunicodechar{Δ}{\ensuremath{\Delta}}
\newunicodechar{∂}{\ensuremath{\partial}}
\newunicodechar{∀}{\ensuremath{\forall}}
\newunicodechar{∃}{\ensuremath{\exists}}
\newunicodechar{√}{\ensuremath{\surd}}
\newunicodechar{−}{\ensuremath{-}}
\newunicodechar{μ}{\ensuremath{\mu}}
\newunicodechar{ρ}{\ensuremath{\rho}}
\newunicodechar{τ}{\ensuremath{\tau}}
\newunicodechar{ω}{\ensuremath{\omega}}
\newunicodechar{Ω}{\ensuremath{\Omega}}
\newunicodechar{χ}{\ensuremath{\chi}}
\newunicodechar{ψ}{\ensuremath{\psi}}
\newunicodechar{Γ}{\ensuremath{\Gamma}}
\newunicodechar{Φ}{\ensuremath{\Phi}}
\newunicodechar{Λ}{\ensuremath{\Lambda}}
\newunicodechar{ℝ}{\ensuremath{\mathbb{R}}}
\newunicodechar{ℤ}{\ensuremath{\mathbb{Z}}}
\newunicodechar{ℕ}{\ensuremath{\mathbb{N}}}
\newunicodechar{△}{\ensuremath{\triangle}}
\newunicodechar{⊗}{\ensuremath{\otimes}}
\begin{document}
\section{2015年杭州电子科技大学601数学分析考研真题}\label{ux5e74ux676dux5ddeux7535ux5b50ux79d1ux6280ux5927ux5b66601ux6570ux5b66ux5206ux6790ux8003ux7814ux771fux9898}

来源：\texttt{HDU/601数学分析/2015/2015年杭州电子科技大学601数学分析考研真题.pdf}。原
PDF 页序：1, 2。

\section{杭州电子科技大学}\label{ux676dux5ddeux7535ux5b50ux79d1ux6280ux5927ux5b66}

\section{2015年攻读硕士学位研究生招生考试}\label{ux5e74ux653bux8bfbux7855ux58ebux5b66ux4f4dux7814ux7a76ux751fux62dbux751fux8003ux8bd5}

\section{《数学分析》试题}\label{ux6570ux5b66ux5206ux6790ux8bd5ux9898}

\begin{quote}
助手整理：按同题答案恢复分段函数及偏导符号。原题三（3）的实值定义域问题、四（5）的严格不等号问题均在整理后的参考解析中给出反例或适用条件；此处保留原命题。
\end{quote}

（试题共4大题，共2页，总分150 分）

姓名 报考专业 准考证号

【所有答案必须写在答题纸上，做在试卷或草稿纸上无效！】

\subsection{一、填空题（本大题共8小题，每小题5分，本大题共40分）}\label{ux4e00ux586bux7a7aux9898ux672cux5927ux9898ux51718ux5c0fux9898ux6bcfux5c0fux98985ux5206ux672cux5927ux9898ux517140ux5206}

\(\lim _ { x \to \infty } \left( { \frac { x + 3 } { x + 1 } } \right) ^ { 4 x + 4 }\)
的值等于

\begin{enumerate}
\def\labelenumi{\arabic{enumi}.}
\setcounter{enumi}{1}
\tightlist
\item
  设
  \(f(x)=\begin{cases}\ln(1+x)/x,&x>0,\\a,&x=0,\\[\sqrt{1+x}-\sqrt{1-x}]/x,&-1<x<0,\end{cases}\)，在
  \(x=0\) 处连续，则必有
\end{enumerate}

a=

\begin{enumerate}
\def\labelenumi{\arabic{enumi}.}
\setcounter{enumi}{2}
\tightlist
\item
  设 \(e^x-e^y=\sin(xy)\)，则 \(y'|_{x=0}=\)\_\_\_\_。
\end{enumerate}

4.设 \(\int f ( x ) d x = \sin { x } + c\) ，则
\(\int f ^ { ( n ) } ( x ) d x =\)

5.设
\(\int _ { 0 } ^ { x } f \left( t \right) d t = \frac { 1 } { 2 } f ( x ) - \frac { 1 } { 2 }\)
，且 \(f ( 0 ) = 1\) ，则f(x)=

6.交换积分次序
\(\int _ { 0 } ^ { 1 } d y \int _ { \sqrt { y } } ^ { \sqrt { 2 - y ^ { 2 } } } f ( x , y ) d x =\)

7.已知 \(z = \arctan \frac { x + y } { x - y }\) ，求dz=

8.幂级数
\(\sum _ { n = 1 } ^ { \infty } { \frac { 5 ^ { n } + ( - 3 ) ^ { n } } { n } } x ^ { n }\)
的收敛半径R=

\subsection{二、计算题（本大题共4小题，本大题共30分）}\label{ux4e8cux8ba1ux7b97ux9898ux672cux5927ux9898ux51714ux5c0fux9898ux672cux5927ux9898ux517130ux5206}

1．（本小题5分）求极限
\(\lim _ { x \to 0 } { \frac { \tan x - x } { x - \sin x } }\)

2.（本小题5分）计算
\(\int _ { 0 } ^ { a } x ^ { 2 n } { \sqrt { a ^ { 2 } - x ^ { 2 } } } d x \left( a > 0 \right)\)
，n为正整数.

3．（本小题10分）设 \(z=x^3f(xy^2,\sin(xy))\)，其中 \(f\)
具有二阶连续偏导数，求 \(\dfrac{\partial^2z}{\partial x\partial y}\)。

4．（本小题10分）计算
\(I = \int _ { \Sigma } \int _ { } ( 2 x + z ) d y d z + z d x d y\)
，其中 \(\Sigma\) 是 \(z = x ^ { 2 } + y ^ { 2 }\)
\(( 0 \leq z \leq 1 )\) ，其曲面正侧法向量与z轴正向的夹角为锐角.

\subsection{三、综合题（本大题共3小题，每小题10分，本大题共30分）}\label{ux4e09ux7efcux5408ux9898ux672cux5927ux9898ux51713ux5c0fux9898ux6bcfux5c0fux989810ux5206ux672cux5927ux9898ux517130ux5206}

1．证明积分 \(\int _ { 0 } ^ { + \infty } x e ^ { - x y } d y\) 在
\([ a , b ] ( a > 0 )\) 上一致收敛，而在{[}0,b{]}上非一致收敛.

2．设曲面
\(S : \frac { x ^ { 2 } } { 2 } + y ^ { 2 } + \frac { z ^ { 2 } } { 4 } = 1\)
，平面 \(\pi : 2 x + 2 y + z + 5 = 0\)

(1)试在曲面S上求平行于平面π的切平面方程.

(2)试求曲面S与平面π之间的最短距离.

\begin{enumerate}
\def\labelenumi{\arabic{enumi}.}
\setcounter{enumi}{2}
\tightlist
\item
  设
  \(f(x,y)=\begin{cases}\dfrac{x^\alpha}{x^2+y^2},&x^2+y^2\neq0,\\0,&x^2+y^2=0,\end{cases}\)，\(\alpha>0\)，讨论在原点的连续性。
\end{enumerate}

\subsection{四、证明题（本大题共5小题，每小题10分，本大题共50分)}\label{ux56dbux8bc1ux660eux9898ux672cux5927ux9898ux51715ux5c0fux9898ux6bcfux5c0fux989810ux5206ux672cux5927ux9898ux517150ux5206}

1.设x\textgreater0，常数a\textgreater e．证明
\(( a + x ) ^ { a } < a ^ { a + x }\)

2．试证椭圆抛物面 \(x ^ { 2 } + y ^ { 2 } = 2 a z\) 和双曲抛物面
\(x ^ { 2 } - y ^ { 2 } = 2 a z\)
被圆柱面\(x ^ { 2 } + y ^ { 2 } = R ^ { 2 }\) 所载得的部分面积相等.(a≠0)

3．设f(x)在{[}a,b{]}上连续，在(a,b)内可导，且
\(f ( x ) \neq 0 , x \in ( a , b )\) ，若\(f ( a ) = f ( b ) = 0\)
，证明对任意实数K存在点 \(\xi ( a < \xi < b )\) ，使
\({ \frac { f ^ { \prime } ( \xi ) } { f ( \xi ) } } = K\)

4.证明 \(f ( x ) = x ^ { 3 }\)
在任意有限闭区间{[}a,b{]}上一致连续，但在{[}0,+∞)非一致连续.

5.设f(x)在{[}0,1{]}上有连续导数且f(O)=0.证明：

\[
{ \Big ( } \operatorname* { m a x } _ { 0 \leq x \leq 1 } f ( x ) { \Big ) } ^ { 2 } < \int _ { 0 } ^ { 1 } { \big [ } f ^ { \prime } ( x ) { \big ] } ^ { 2 } d x .
\]
\end{document}
